\begin{textsample}{POS Dim 9 – global\_south\_2025\_06 – Score 1.00 – global\_south\_2025\_06/125536427.txt}  \label{ex:f9_pos_177}
US vetoes UN Security Council demand for Gaza ceasefire June 5, 2025 4:57 pm The United States vetoed a draft UN Security Council resolution that demanded an"immediate, unconditional and permanent ceasefire"between Israel and Hamas militants in Gaza.Credit : UN Photo/ Mark Garten UNITED NATIONS/CAIRO/JERUSALEM -- The United States on Wednesday vetoed a draft United Nations ( UN ) Security Council resolution that demanded an"immediate, unconditional and permanent ceasefire"between Israel and Hamas militants in Gaza and \textbf{unhindered} aid access across the war-torn enclave.
The other 14 countries on the council voted in favor of the draft as a humanitarian crisis grips the enclave of more than 2 million people, where famine looms and aid has only trickled in since Israel lifted an 11-week blockade last month."The United States has been clear : We would not support any measure that fails to condemn Hamas and does not call for Hamas to disarm and leave Gaza,"Acting US Ambassador to the UN Dorothy Shea told the council before the @ @ @ @ @ @ @ @ @ @ to broker a ceasefire.
Washington is Israel ' s biggest ally and arms supplier.
The Security Council vote came as Israel pushes ahead with an offensive in Gaza after ending a two-month truce in March.
Gaza health authorities said Israeli strikes killed 45 people on Wednesday, while Israel said a soldier died in fighting.
Israel has rejected calls for an unconditional or permanent ceasefire, saying Hamas can not stay in Gaza.
Israel ' s UN Ambassador Danny Danon told the council members who voted in favor of the draft :"You chose appeasement and submission.
You chose a road that does not lead to peace.
Only to more terror."Hamas condemned the US veto, describing it as showing"the US administration ' s blind bias"towards Israel.
The draft Security Council resolution had also demanded the immediate and unconditional release of all hostages held by Hamas and others.
RIVAL AID OPERATIONS The war in Gaza has raged since 2023 after Hamas militants killed 1,200 people in Israel in an October 7 @ @ @ @ @ @ @ @ @ @, according to Israeli tallies.
Many of those killed or captured were civilians.
Israel responded with a military campaign that has killed over 54,000 Palestinians, according to Gaza health authorities.
They say civilians have borne the brunt of the attacks and that thousands more bodies have been lost under rubble.
Under global pressure, Israel allowed limited UN-led deliveries to resume on May 19.
A week later a controversial new aid distribution system was launched by the Gaza Humanitarian Foundation ( GHF ), backed by the US and Israel.
Israel has long accused Hamas of stealing aid, which the group denies.
Israel and the US are urging the UN to work through the GHF, which is using private US security and logistics companies to transport aid into Gaza for distribution at so-called secure distribution sites."No one wants to see Palestinian civilians in Gaza go hungry or thirsty,"Ms.
Shea told the Security Council, adding that the draft resolution did not"acknowledge the disastrous shortcomings of the prior method @ @ @ @ @ @ @ @ @ @ aid groups have refused to work with the GHF because they say it is not neutral, militarizes aid and forces the displacement of Palestinians.
No aid was distributed by the US-backed GHF on Wednesday as it pressed the Israeli military to boost civilian safety beyond the perimeter of its so-called secure distribution sites after a deadly incident on Tuesday.
The GHF said it has asked the Israeli military to"guide foot traffic in a way that minimizes confusion or escalation risks"near military positions, provide clearer civilian guidance and enhance training of soldiers on civilian safety. ' DELAYS AND DENIALS ' The GHF posted on Facebook that"ongoing maintenance work"would delay the opening of its distribution sites on Thursday.
It said on Tuesday that it has so far distributed more than seven million meals since it started operations.
Despite US and Israeli criticism of the UN-led Gaza aid operation, a US ceasefire plan proposes the delivery of aid by the UN, the Red Crescent and other agreed channels.
Israel has agreed @ @ @ @ @ @ @ @ @ @ the US has rejected as"totally unacceptable."Ahead of the UN Security Council vote, UN Aid Chief Tom Fletcher again appealed for the UN and aid groups to be allowed to assist people in Gaza, stressing that they have a plan, supplies and experience."Open the crossings -- all of them.
Let in lifesaving aid at scale, from all directions.
Lift the restrictions on what and how much aid we can bring in.
Ensure our convoys aren't held up by delays and denials,"Mr.
Fletcher said in a statement.
The UN has long-blamed Israel and lawlessness in the enclave for hindering the delivery of aid into Gaza and its distribution throughout the war zone."Enough of suffering of civilians.
Enough of food being used as a weapon.
Enough is enough is enough,"Slovenia ' s UN Ambassador Samuel Zbogar told the Security Council.
A similar humanitarian-focused draft resolution is now expected to be put to a vote in the 193-member UN General Assembly @ @ @ @ @ @ @ @ @ @ would likely pass, diplomats said.
Mr.
Danon warned :"Don't waste more of your time, because no resolution, no vote, no moral failure, will stand in our way."-- Reuters

% matched lemmas: unhindered
\end{textsample}
